\chapter{Funções Usuais}\label{ch:b1:functions}

A análise só é tão útil quanto o acervo de funções que se domina. À
coleção herdada do ensino médio --- potências, exponencial,
logaritmo, funções trigonométricas --- este capítulo acrescenta as suas
funções inversas ($\arcsin$, $\arccos$, $\arctan$) e a família
hiperbólica. As derivadas são usadas livremente no nível do ensino
médio; a teoria por trás das funções inversas é completada no
\cref{ch:b1:continuity,ch:b1:derivative}.

\section{Exponencial, logaritmo, potências}

\begin{proposition}[Retomada e caracterização]\label{prop:b1:functions:expln}
$\exp \colon \R \to \intoo{0}{+\infty}$ e $\ln \colon
\intoo{0}{+\infty} \to \R$ são bijeções recíprocas, estritamente
crescentes, com
\[
\exp(x + y) = \exp x \exp y, \qquad \ln(xy) = \ln x + \ln y,
\qquad \exp' = \exp, \qquad \ln'(x) = \frac 1x .
\]
Além disso, $\exp$ é a \emph{única} função derivável $f \colon
\R \to \R$ com $f' = f$ e $f(0) = 1$.
\end{proposition}

\begin{proof}
A retomada é matéria do ensino médio. Para a unicidade, sejam $f' = f$,
$f(0) = 1$, e ponha $g(x) = f(x)\,\eu^{-x}$. Então $g' = f'\eu^{-x} -
f\eu^{-x} = 0$, de modo que $g$ é constante igual a $g(0) = 1$: $f = \exp$.
\end{proof}

\begin{definition}[Potências gerais]\label{def:b1:functions:powers}
Para $x > 0$ e $\alpha \in \R$:
$\;x^\alpha = \eu^{\alpha \ln x}$\index{função potência}. Para $a > 0$,
$a \neq 1$, o \emph{logaritmo na base $a$} é $\log_a x =
\frac{\ln x}{\ln a}$, a inversa de $x \mapsto a^x$.
\end{definition}

\begin{example}[Resolvendo equações exponenciais]\label{ex:b1:functions:expequations}
Resolva $2^x = 5^{\,x-1}$ em $\R$. Os dois lados são positivos, então tome
logaritmos --- um passo reversível:
\[
x\ln2 = (x - 1)\ln5
\iff x(\ln2 - \ln5) = -\ln5
\iff x = \frac{\ln5}{\ln5 - \ln2} = \frac{\ln 5}{\ln\frac52}
\approx 1.756 .
\]
Em seguida, resolva $x^{\sqrt2} = 3$ para $x > 0$: eleve à potência
$\frac1{\sqrt2}$ (isto é, aplique a bijeção recíproca):
$x = 3^{1/\sqrt2} = \eu^{(\ln 3)/\sqrt2} \approx 2.175$. A ideia:
toda equação que mistura potências se desenrola por meio de $\ln$ e
$\exp$, porque a definição $x^\alpha = \eu^{\alpha\ln x}$
reduz toda manipulação de potências à aritmética dos expoentes ---
mas apenas no domínio $x > 0$ em que essa definição vive.
\end{example}

\begin{example}[Tempos de duplicação]\label{ex:b1:functions:doubling}
Uma grandeza cresce $3\%$ a cada passo: após $n$ passos ela fica
multiplicada por $(1.03)^n$. Quando ela dobra? Resolva $(1.03)^n
\geq 2$:
\[
n\ln(1.03) \geq \ln2
\iff n \geq \frac{\ln 2}{\ln 1.03}
= \frac{0.6931}{0.02956} \approx 23.4 ,
\]
de modo que a primeira duplicação ocorre no passo $24$. (A
``regra dos $72$'' dos financistas, que estima o tempo de duplicação como
$72$ dividido pela taxa em porcentagem, é este cálculo com
a aproximação $\ln(1 + x) \approx x$, quantificada no
\cref{ch:b1:taylor}.) Os processos
exponenciais são mais bem tratados através dos seus logaritmos: nessa
escala, o crescimento é linear e as perguntas se tornam divisões.
\end{example}

\begin{example}[Qual é o tamanho de $2^{2026}$?]\label{ex:b1:functions:digits}
O número de algarismos decimais de um inteiro $N \geq 1$ é
$\floor{\log_{10} N} + 1$ (com efeito, $N$ tem $d$ algarismos exatamente quando
$10^{d-1} \leq N < 10^d$, isto é, $d - 1 \leq \log_{10}N < d$). Para
$N = 2^{2026}$:
\[
\log_{10} 2^{2026} = 2026\,\log_{10}2
= 2026 \times 0.301030 = 609.887 ,
\]
de modo que $2^{2026}$ tem $610$ algarismos. A parte fracionária traz um
bônus: $10^{0.887} \approx 7.7$, de modo que o número \emph{começa} por
$7$. Uma multiplicação respondeu a uma pergunta sobre um número que
ninguém jamais escreverá por extenso --- os logaritmos comprimem
tamanho multiplicativo em tamanho aditivo, o que é toda a sua razão
histórica de ser (\cref{ex:b1:structures:expmorphism} torna essa frase precisa).
\end{example}

\begin{proposition}[Regras de potências e comparação de crescimentos]\label{prop:b1:functions:powerrules}
Para $x, y > 0$ e $\alpha, \beta \in \R$:
\[
x^{\alpha+\beta} = x^\alpha x^\beta, \quad
(x^\alpha)^\beta = x^{\alpha\beta}, \quad
(xy)^\alpha = x^\alpha y^\alpha, \quad
(x^\alpha)' = \alpha\, x^{\alpha - 1}.
\]
Escala de crescimento\index{comparação de crescimentos} quando $x \to +\infty$, para todos $\alpha > 0$ e $\beta > 0$:
\[
\frac{(\ln x)^\beta}{x^\alpha} \longrightarrow 0,
\qquad
\frac{x^\alpha}{\eu^{\beta x}} \longrightarrow 0 .
\]
\end{proposition}

\begin{proof}
As identidades transcrevem as de $\exp$ e $\ln$ através da
definição. Em detalhe, para a segunda (a menos evidente): $x^\alpha
> 0$ e $\ln(x^\alpha) = \alpha\ln x$ (aplique $\ln$ à
definição), de modo que
\[
(x^\alpha)^\beta = \eu^{\beta\ln(x^\alpha)}
= \eu^{\beta\alpha\ln x} = x^{\alpha\beta} ;
\]
a primeira e a terceira são os mesmos desenrolamentos de duas linhas via
$\exp(u + v) = \exp u\exp v$ e $\ln(xy) = \ln x + \ln y$. A
derivada é a regra da cadeia: $(\eu^{\alpha\ln x})' =
\frac{\alpha}{x} \eu^{\alpha\ln x}$.

Comparações: de $\frac{\ln t}{t} \to 0$ (demonstrado no volume do
ensino médio), substitua $t = x^{\alpha/\beta}$:
$\frac{\ln x}{x^{\alpha/\beta}} \to 0$, e depois eleve à potência
$\beta$. Para a segunda, substitua $x = \ln u$ na primeira.
\end{proof}

\begin{example}[Dois limites que todo leitor deve dominar]\label{ex:b1:functions:xx}
Quanto valem $\lim_{x \to 0^+} x^x$ e $\lim_{x \to +\infty}
x^{1/x}$? As duas potências estão definidas via a exponencial, de modo
que é preciso resolver antes o expoente:
\[
x^x = \eu^{x\ln x}, \qquad
x\ln x = -\frac{\ln(1/x)}{1/x} \longrightarrow 0
\quad (x \to 0^+),
\qquad\text{logo } x^x \longrightarrow \eu^0 = 1 ;
\]
e $x^{1/x} = \eu^{(\ln x)/x} \to \eu^0 = 1$ quando $x \to +\infty$,
diretamente pela \omterm{prop:b1:functions:powerrules}{comparação de crescimentos}. A ideia: uma potência
indeterminada (do tipo $0^0$ ou $\infty^0$) é \emph{sempre} tratada
reescrevendo-se $u^v = \eu^{v\ln u}$ e analisando o produto $v \ln u$
--- nunca adivinhando separadamente a partir da base e do expoente. A
função $x \mapsto \frac{\ln x}{x}$ que decidiu os dois limites é
estudada exaustivamente no problema de fim de semana deste capítulo.
\end{example}

\begin{omfigure}
\begin{tikzpicture}
\begin{axis}[omaxis, width=9.5cm, height=6cm,
    xmin=0, xmax=5.2, ymin=-1.8, ymax=8.5,
    xtick={1,2,3,4,5}, ytick={2,4,6,8}]
  \addplot[omMeth, thick, domain=0.02:5, samples=200] {ln(x)}
    node[pos=0.97, below, font=\small] {$\ln x$};
  \addplot[omDef, thick, domain=0:5, samples=100] {sqrt(x)}
    node[pos=0.9, above, font=\small] {$\sqrt x$};
  \addplot[omProp, thick, domain=0:5] {x}
    node[pos=0.72, left, font=\small] {$x$};
  \addplot[omThm, thick, domain=0:2.14, samples=120] {exp(x)}
    node[pos=0.85, left, font=\small] {$\eu^x$};
\end{axis}
\end{tikzpicture}

{\small A escala de crescimento da
\cref{prop:b1:functions:powerrules}: perto da borda direita, $\ln x$
mal passou de $1.6$, enquanto $\eu^x$ já saiu do quadro. Toda
razão (logaritmo sobre potência, potência sobre exponencial) tende a $0$
--- a figura apenas sugere o que as substituições da demonstração tornam
exato.}
\end{omfigure}

\begin{remark}[Armadilhas frequentes com potências e logaritmos]\label{rem:b1:functions:pitfalls}
\begin{enumerate}
  \item \emph{Domínios.} $x^\alpha$ para $\alpha$ irracional exige
        $x > 0$; $(-8)^{1/3}$ é melhor evitado em favor de ``a
        raiz cúbica real de $-8$'', porque a regra
        $(x^\alpha)^\beta = x^{\alpha\beta}$ falha silenciosamente em
        números negativos: $\bigl((-8)^2\bigr)^{1/6} = 2 \neq -2$.
  \item \emph{$\sqrt{x^2} = \abs x$}, e não $x$: esquecer o
        valor absoluto é a fonte clássica de soluções negativas
        perdidas.
  \item \emph{$\ln(xy) = \ln x + \ln y$ exige $x, y > 0$.} Para
        argumentos negativos, $\ln(xy)$ pode estar definido enquanto o
        lado direito não está.
  \item \emph{Qual função cresce?} Em $x^\alpha$ varia a
        \emph{base}; em $a^x$ varia o \emph{expoente}.
        Seus crescimentos são radicalmente diferentes
        (\cref{prop:b1:functions:powerrules}), e expressões
        híbridas como $x^{\ln x}$ ou $x^{1/x}$ devem ser
        reescritas como $\eu^{(\cdot)}$ antes de qualquer raciocínio ---
        como no \cref{ex:b1:functions:xx}.
\end{enumerate}
\end{remark}

\section{Funções trigonométricas inversas}

\begin{definition}[$\arcsin$, $\arccos$, $\arctan$]\label{def:b1:functions:arc}
As restrições
\[
\sin \colon \intcc{-\tfrac\pi2}{\tfrac\pi2} \to \intcc{-1}{1},
\qquad
\cos \colon \intcc{0}{\pi} \to \intcc{-1}{1},
\qquad
\tan \colon \intoo{-\tfrac\pi2}{\tfrac\pi2} \to \R
\]
são bijeções estritamente monótonas. Suas aplicações inversas escrevem-se
$\arcsin$\index{arco-seno}, $\arccos$\index{arco-cosseno} e
$\arctan$\index{arco-tangente}. Assim, por exemplo, $y = \arcsin x$
($x \in \intcc{-1}{1}$) é \emph{o} ângulo de
$\intcc{-\frac\pi2}{\frac\pi2}$ cujo seno é $x$.
\end{definition}

\begin{omfigure}
\begin{tikzpicture}
\begin{axis}[omaxis, width=5.6cm, height=5.6cm,
    xmin=-1.25, xmax=1.25, ymin=-1.8, ymax=1.8,
    xtick={-1,1}, ytick={-1.5708,1.5708},
    yticklabels={$-\frac{\pi}{2}$,$\frac{\pi}{2}$}]
  \addplot[omDef, very thick, domain=-1:1, samples=200]
    {rad(asin(x))};
\end{axis}
\begin{scope}[xshift=5.3cm]
\begin{axis}[omaxis, width=5.6cm, height=5.6cm,
    xmin=-1.25, xmax=1.25, ymin=-0.4, ymax=3.4,
    xtick={-1,1}, ytick={1.5708,3.1416},
    yticklabels={$\frac{\pi}{2}$,$\pi$}]
  \addplot[omThm, very thick, domain=-1:1, samples=200]
    {rad(acos(x))};
\end{axis}
\end{scope}
\begin{scope}[xshift=10.6cm]
\begin{axis}[omaxis, width=5.6cm, height=5.6cm,
    xmin=-4, xmax=4, ymin=-1.9, ymax=1.9,
    xtick={-2,2}, ytick={-1.5708,1.5708},
    yticklabels={$-\frac{\pi}{2}$,$\frac{\pi}{2}$}]
  \addplot[omMeth, very thick, domain=-4:4, samples=200]
    {rad(atan(x))};
  \addplot[dashed, gray, domain=-4:4] {1.5708};
  \addplot[dashed, gray, domain=-4:4] {-1.5708};
\end{axis}
\end{scope}
\end{tikzpicture}

{\small Da esquerda para a direita: $\arcsin$, $\arccos$, $\arctan$. Cada
uma herda o seu gráfico da função direta restrita, por reflexão
na reta $y = x$.}
\end{omfigure}

\begin{proposition}[Derivadas]\label{prop:b1:functions:arcderiv}
No interior dos seus domínios:
\[
\arcsin' x = \frac{1}{\sqrt{1 - x^2}},
\qquad
\arccos' x = \frac{-1}{\sqrt{1 - x^2}},
\qquad
\arctan' x = \frac{1}{1 + x^2}.
\]
\end{proposition}

\begin{proof}
Antecipando a regra de derivação da função inversa, demonstrada no
\cref{ch:b1:derivative}: se $f$ é uma bijeção derivável e
$f' \neq 0$, então $(f^{-1})'(x) = \frac{1}{f'(f^{-1}(x))}$. Para
$\arcsin$: com $y = \arcsin x$,
\[
\arcsin' x = \frac{1}{\cos y}
= \frac{1}{\sqrt{1 - \sin^2 y}} = \frac{1}{\sqrt{1 - x^2}},
\]
em que $\cos y = +\sqrt{1 - \sin^2 y}$, pois $y \in
\intoo{-\frac\pi2}{\frac\pi2}$ força $\cos y > 0$. Analogamente,
$\arccos' x = \frac{1}{-\sin y}$ com $\sin y > 0$ em
$\intoo{0}{\pi}$, e $\arctan' x = \frac{1}{1 + \tan^2 y} =
\frac{1}{1 + x^2}$, usando $\tan' = 1 + \tan^2$.
\end{proof}

\begin{example}[Reconhecendo uma constante escondida]\label{ex:b1:functions:hiddenconst}
Estude $g(x) = \arctan\dfrac{1 - x}{1 + x}$ em $\intoo{-1}
{+\infty}$. Regra da cadeia e um cálculo curto:
\[
g'(x) = \frac{1}{1 + \bigl(\frac{1-x}{1+x}\bigr)^2}\cdot
\frac{-(1+x) - (1-x)}{(1+x)^2}
= \frac{-2}{(1+x)^2 + (1-x)^2}
= \frac{-1}{1 + x^2} ,
\]
pois $(1+x)^2 + (1-x)^2 = 2 + 2x^2$. Logo $g' = -\arctan'$: a
função $g + \arctan$ tem derivada nula no \emph{intervalo}
$\intoo{-1}{+\infty}$ e, portanto, é constante nele; o seu valor em $x =
0$ é $\arctan 1 + \arctan 0 = \frac\pi4$. Conclusão:
\[
\arctan\frac{1 - x}{1 + x} = \frac\pi4 - \arctan x
\qquad (x > -1) .
\]
Em $\intoo{-\infty}{-1}$ o mesmo cálculo de derivada vale, mas
a constante é diferente ($-\frac{3\pi}4$: avalie o limite quando
$x \to -\infty$). A ideia: ``derivada nula implica
constante'' é um enunciado intervalo por intervalo --- exatamente a
sutileza explorada no \cref{exo:b1:functions:6} e no
\cref{exo:b1:functions:10}.
\end{example}

\begin{proposition}[Identidades usuais]\label{prop:b1:functions:arcidentities}
\begin{enumerate}
  \item Para $x \in \intcc{-1}{1}$: $\arcsin x + \arccos x =
        \dfrac{\pi}{2}$.
  \item Para $x > 0$: $\arctan x + \arctan\dfrac 1x = \dfrac{\pi}{2}$
        (e $-\frac\pi2$ para $x < 0$).
  \item $\sin(\arccos x) = \cos(\arcsin x) = \sqrt{1 - x^2}$;
        $\;\tan(\arcsin x) = \dfrac{x}{\sqrt{1 - x^2}}$ para
        $\abs x < 1$.
\end{enumerate}
\end{proposition}

\begin{proof}
(1) A derivada de $x \mapsto \arcsin x + \arccos x$ é nula em
$\intoo{-1}{1}$ (\cref{prop:b1:functions:arcderiv}), de modo que a função
é constante ali, igual ao seu valor $\frac\pi2$ em $0$; os valores nos
extremos $\pm 1$ verificam-se diretamente ($\frac\pi2 + 0$ e $-\frac\pi2 +
\pi$).

(2) Mesmo método em $\intoo{0}{+\infty}$: a derivada é
$\frac{1}{1+x^2} + \frac{-1/x^2}{1 + 1/x^2} = \frac{1}{1+x^2} -
\frac{1}{x^2 + 1} = 0$, e em $x = 1$ a soma vale $2\arctan 1 =
\frac\pi2$. Para $x < 0$, use a imparidade de $\arctan$.

(3) Com $y = \arccos x \in \intcc{0}{\pi}$: $\sin y \geq 0$, logo
$\sin y = \sqrt{1 - \cos^2 y} = \sqrt{1 - x^2}$; do mesmo modo para
$\cos(\arcsin x)$, e a fórmula da tangente é o quociente.
\end{proof}

\begin{remark}
$\arcsin(\sin\theta) = \theta$ vale \emph{apenas} para $\theta \in
\intcc{-\frac\pi2}{\frac\pi2}$: por exemplo,
$\arcsin(\sin\pi) = 0 \neq \pi$. A composição no outro
sentido, $\sin(\arcsin x) = x$, é válida em todo
$\intcc{-1}{1}$.
\end{remark}

\begin{example}[Voltando ao intervalo principal]
\label{ex:b1:functions:wrap}
Calcule
\[
\arctan\Bigl(\tan\frac{3\pi}4\Bigr)
\qquad\text{e}\qquad
\arctan\Bigl(\tan\frac{17\pi}5\Bigr).
\]
A receita: substitua o ângulo pelo único ângulo de
$\intoo{-\frac\pi2}{\frac\pi2}$ com a mesma tangente, isto é, subtraia o
múltiplo adequado de $\pi$ (o período de $\tan$). Primeiro: $\frac{3\pi}4 - \pi =
-\frac\pi4$, de modo que a resposta é $-\frac\pi4$. Segundo: $\frac{17\pi}
5 - 3\pi = \frac{2\pi}5 \in \intoo{-\frac\pi2}{\frac\pi2}$, de modo que
a resposta é $\frac{2\pi}5$. O cálculo é uma divisão euclidiana
do ângulo por $\pi$ disfarçada --- e as receitas análogas
para $\arcsin$ (refletir para
$\intcc{-\frac\pi2}{\frac\pi2}$, período $2\pi$) e $\arccos$
(refletir para $\intcc0\pi$) conduzem o \cref{exo:b1:functions:1} e
a resposta por partes do \cref{exo:b1:functions:10}.
\end{example}

\begin{example}[Somando arco-tangentes com segurança]\label{ex:b1:functions:arctansum}
Demonstremos que
\[
\arctan\frac12 + \arctan\frac15 + \arctan\frac18 = \frac\pi4 .
\]
Dois ingredientes: a fórmula da tangente da soma e --- o passo que os
iniciantes esquecem --- uma \emph{localização} da soma. Primeiro,
$\tan(u + v) = \frac{\tan u + \tan v}{1 - \tan u\tan v}$ com
$u = \arctan\frac12$, $v = \arctan\frac15$ dá
\[
\tan(u + v) = \frac{\frac12 + \frac15}{1 - \frac1{10}}
= \frac{7/10}{9/10} = \frac79,
\qquad\text{logo}\qquad
\tan\Bigl(u + v + \arctan\frac18\Bigr)
= \frac{\frac79 + \frac18}{1 - \frac7{72}}
= \frac{65/72}{65/72} = 1 .
\]
Segundo, cada um dos três ângulos está em $\intoo0{\frac\pi4}$
(seus argumentos são menores que $1$), de modo que a soma está em
$\intoo0{\frac{3\pi}4}$; o único ângulo ali com tangente $1$ é
$\frac\pi4$. Sem a localização, a conclusão seria
``$\frac\pi4$ a menos de um múltiplo de $\pi$'' --- meia demonstração. A
mesma disciplina em dois passos conduz o \cref{exo:b1:functions:8} e o
\cref{exo:b1:functions:12}.
\end{example}

\begin{example}[Uma equação com arco-tangente]\label{ex:b1:functions:arctaneq}
Resolva $\arctan x + \arctan 2x = \dfrac\pi4$. Primeiro localize: o
lado esquerdo tem o sinal de $x$ (os dois termos têm), de modo que toda
solução tem $x > 0$, e então a soma está em $\intoo0\pi$. Tome tangentes
(\omterm{def:b1:logic:inj}{injetiva} em nenhum intervalo de comprimento $\pi$, mas combinada com a
localização isso bastará): a fórmula da soma dá
\[
\tan\bigl(\arctan x + \arctan 2x\bigr)
= \frac{3x}{1 - 2x^2} = 1
\iff 2x^2 + 3x - 1 = 0
\iff x = \frac{-3 \pm \sqrt{17}}4 .
\]
A raiz negativa é descartada pela localização; para o
candidato positivo $x_0 = \frac{\sqrt{17} - 3}4 \approx 0.28$,
a soma está em $\intoo0\pi$ com tangente $1$, e o único ângulo
desses é $\frac\pi4$ (em $\intoo{\frac\pi2}\pi$ a tangente é
negativa): $x_0$ é a única solução. Note a forma do
argumento: tomar tangentes pode criar soluções, nunca perdê-las,
de modo que se resolve a equação polinomial \emph{e depois} se filtra
pela localização --- a mesma disciplina de verificação a posteriori que
se usa ao elevar uma equação ao quadrado.
\end{example}

\section{Funções hiperbólicas}

\begin{definition}[$\cosh$, $\sinh$, $\tanh$]\label{def:b1:functions:hyperbolic}
Para $x \in \R$:
\[
\cosh x = \frac{\eu^x + \eu^{-x}}{2},
\qquad
\sinh x = \frac{\eu^x - \eu^{-x}}{2},
\qquad
\tanh x = \frac{\sinh x}{\cosh x}
\]
(\emph{cosseno, seno e tangente hiperbólicos}\index{funções
hiperbólicas}). $\cosh$ é par; $\sinh$ e $\tanh$ são ímpares.
\end{definition}

\begin{proposition}[Propriedades básicas]\label{prop:b1:functions:hyprules}
Para todos $x, y \in \R$:
\begin{enumerate}
  \item $\cosh^2 x - \sinh^2 x = 1$;
  \item $\cosh' = \sinh$, $\;\sinh' = \cosh$, $\;\tanh' = 1 - \tanh^2
        = \dfrac{1}{\cosh^2}$;
  \item $\cosh(x+y) = \cosh x\cosh y + \sinh x \sinh y$ e
        $\sinh(x+y) = \sinh x \cosh y + \cosh x \sinh y$;
  \item $\sinh$ é uma bijeção estritamente crescente de $\R$ sobre $\R$;
        $\cosh$ restrita a $\R_+$ é uma bijeção estritamente crescente
        sobre $\intco{1}{+\infty}$; $\tanh$ é uma bijeção estritamente
        crescente de $\R$ sobre $\intoo{-1}{1}$.
\end{enumerate}
\end{proposition}

\begin{proof}
(1) $\bigl(\frac{\eu^x + \eu^{-x}}{2}\bigr)^2 - \bigl(\frac{\eu^x -
\eu^{-x}}{2}\bigr)^2 = \frac{(\eu^{2x} + 2 + \eu^{-2x}) - (\eu^{2x} -
2 + \eu^{-2x})}{4} = 1$.

(2) Derive as fórmulas de definição; para $\tanh$, a regra do
quociente dá $\frac{\cosh^2 - \sinh^2}{\cosh^2}$, que é ao mesmo tempo
$1 - \tanh^2$ e $\frac{1}{\cosh^2}$ por (1).

(3) Expanda os lados direitos usando as definições. Em detalhe,
para a primeira fórmula:
\[
\cosh x\cosh y + \sinh x\sinh y
= \frac{(\eu^x + \eu^{-x})(\eu^y + \eu^{-y})
+ (\eu^x - \eu^{-x})(\eu^y - \eu^{-y})}{4} ;
\]
dos oito produtos, os quatro ``mistos''
($\eu^{x-y}$, $\eu^{y-x}$) se cancelam aos pares, ao passo que
$\eu^{x+y}$ e $\eu^{-(x+y)}$ aparecem cada um duas vezes:
o total é $\frac{2\eu^{x+y} + 2\eu^{-(x+y)}}{4} = \cosh(x+y)$.
A fórmula do seno é idêntica, com os termos mistos sobrevivendo
em lugar dos outros.

(4) $\sinh' = \cosh \geq 1 > 0$, de modo que $\sinh$ é estritamente crescente,
com limites $\pm\infty$ (termo dominante $\pm\frac{\eu^{\abs
x}}{2}$);
o enunciado da bijeção decorre então do teorema do valor intermediário
(usado no nível do ensino médio; tratamento sistemático no
\cref{ch:b1:continuity}). Em $\R_+$, $\cosh' = \sinh \geq 0$,
anulando-se apenas em $0$: estritamente crescente de $\cosh 0 = 1$ a
$+\infty$. E $\tanh' > 0$, com limites $\pm 1$ em $\pm\infty$: em
detalhe,
\[
\tanh x = \frac{\eu^x - \eu^{-x}}{\eu^x + \eu^{-x}}
= \frac{1 - \eu^{-2x}}{1 + \eu^{-2x}}
\xrightarrow[x \to +\infty]{} 1 ,
\]
depois de dividir numerador e denominador por $\eu^x$; a imparidade
dá o limite $-1$ em $-\infty$; a função estritamente crescente
$\tanh$ leva, portanto, $\R$ sobre $\intoo{-1}1$.
\end{proof}

\begin{omfigure}
\begin{tikzpicture}
\begin{axis}[omaxis, width=6.4cm, height=6cm,
    xmin=-2.6, xmax=2.6, ymin=-3.6, ymax=3.9,
    xtick={-2,-1,1,2}, ytick={-3,-2,-1,1,2,3}]
  \addplot[omDef, very thick, domain=-2.05:2.05, samples=120]
    {cosh(x)} node[pos=0.1, right, font=\small] {$\cosh$};
  \addplot[omThm, very thick, domain=-2.05:2.05, samples=120]
    {sinh(x)} node[pos=0.93, above left, font=\small] {$\sinh$};
  \addplot[dashed, gray, domain=0:2.5] {exp(x)/2};
\end{axis}
\begin{scope}[xshift=7.4cm]
\begin{axis}[omaxis, width=6.4cm, height=6cm,
    xmin=-3.2, xmax=3.2, ymin=-1.5, ymax=1.5,
    xtick={-2,-1,1,2}, ytick={-1,1}]
  \addplot[omMeth, very thick, domain=-3:3, samples=150]
    {tanh(x)} node[pos=0.9, above left, font=\small] {$\tanh$};
  \addplot[dashed, gray, domain=-3:3] {1};
  \addplot[dashed, gray, domain=-3:3] {-1};
\end{axis}
\end{scope}
\end{tikzpicture}

{\small À esquerda: $\cosh$ e $\sinh$, assintoticamente colados a
$\frac{\eu^x}{2}$ (tracejado). À direita: $\tanh$, crescente de $-1$ a
$1$.}
\end{omfigure}

\begin{remark}[Por que ``hiperbólicas''?]
O ponto $(\cosh t, \sinh t)$ percorre o ramo $x > 0$ da
hipérbole $x^2 - y^2 = 1$ (pela
\cref{prop:b1:functions:hyprules}~(1)), exatamente como $(\cos t, \sin
t)$
percorre o círculo $x^2 + y^2 = 1$. Toda identidade circular tem
uma irmã hiperbólica, com trocas de sinal governadas por
$\sinh^2 \leftrightarrow -\sin^2$.
\end{remark}

\begin{example}[A fórmula de adição para $\tanh$]\label{ex:b1:functions:tanhadd}
Dividindo as duas fórmulas de adição da
\cref{prop:b1:functions:hyprules}~(3) por $\cosh x\cosh y$:
\[
\tanh(x + y)
= \frac{\sinh x\cosh y + \cosh x\sinh y}
{\cosh x\cosh y + \sinh x\sinh y}
= \frac{\tanh x + \tanh y}{1 + \tanh x\,\tanh y} ,
\]
a irmã hiperbólica da fórmula de adição da tangente --- com um
$+$ onde a trigonometria tem um $-$. Um dividendo: como $\abs{\tanh}
< 1$, o lado direito é uma regra de ``adição de velocidades'' que
nunca sai de $\intoo{-1}1$: se $u, v \in \intoo{-1}1$, então
$\frac{u + v}{1 + uv} \in \intoo{-1}1$ também (escreva $u = \tanh
a$, $v = \tanh b$, o que é possível pela bijetividade, e leia a fórmula
ao contrário). Verificar isso algebricamente, sem \omterm{def:b1:functions:hyperbolic}{funções
hiperbólicas}, é um exercício ligeiramente penoso; parametrizar por
$\tanh$ o torna uma linha --- a mesma estratégia que as funções
circulares oferecem para o círculo unitário.
\end{example}

\begin{proposition}[Funções hiperbólicas inversas]\label{prop:b1:functions:invhyp}
As inversas concedidas pela \cref{prop:b1:functions:hyprules}~(4) têm
formas fechadas:
\[
\operatorname{arsinh} x = \ln\bigl(x + \sqrt{x^2 + 1}\bigr)
\ (x \in \R),
\qquad
\operatorname{arcosh} x = \ln\bigl(x + \sqrt{x^2 - 1}\bigr)
\ (x \geq 1),
\]
\[
\operatorname{artanh} x = \frac 12 \ln\frac{1 + x}{1 - x}
\ (\abs x < 1),
\]
com derivadas $\frac{1}{\sqrt{x^2+1}}$, $\frac{1}{\sqrt{x^2-1}}$
($x > 1$) e $\frac{1}{1 - x^2}$, respectivamente.
\end{proposition}

\begin{proof}
Para $\operatorname{arsinh}$: resolva $x = \sinh y = \frac{\eu^y -
\eu^{-y}}{2}$. Pondo $u = \eu^y > 0$: $u^2 - 2xu - 1 = 0$, de modo que $u = x
+ \sqrt{x^2 + 1}$ (a raiz $x - \sqrt{x^2+1}$ é negativa), e $y =
\ln(x + \sqrt{x^2+1})$. As outras duas são cálculos idênticos
($u^2 - 2xu + 1 = 0$ para $\operatorname{arcosh}$, mantendo-se a raiz
$\geq 1$; um rearranjo de duas linhas para $\operatorname{artanh}$).
Derivadas: derive as expressões logarítmicas, por exemplo,
\[
\bigl(\ln(x + \sqrt{x^2+1})\bigr)'
= \frac{1 + \frac{x}{\sqrt{x^2+1}}}{x + \sqrt{x^2+1}}
= \frac{1}{\sqrt{x^2 + 1}} . \qedhere
\]
\end{proof}

\begin{example}[Formas fechadas em ação]\label{ex:b1:functions:arcoshtwo}
A solução de $\cosh t = 2$ com $t \geq 0$ é, pela forma
fechada, $t = \operatorname{arcosh} 2 = \ln(2 + \sqrt3) \approx
1.317$; a outra solução é $-t$, pela paridade --- e, de fato,
$\ln(2 - \sqrt3) = \ln\frac1{2 + \sqrt3} = -\ln(2 + \sqrt3)$: as
duas raízes $u = \eu^t$ da equação do segundo grau $u^2 - 4u + 1 = 0$ são
recíprocas, como o seu produto $1$ (Girard) exige. Esse cálculo
minúsculo exibe o padrão geral: equações hiperbólicas
convertem-se em equações do segundo grau em $\eu^t$, e a simetria $t \mapsto -t$
aparece como a simetria $u \mapsto 1/u$ da equação quadrática ---
vale lembrar disso ao resolver o \cref{exo:b1:functions:7}.
\end{example}

\begin{method}[Escolhendo a forma primitiva certa]\label{met:b1:functions:primitive}
Os três padrões de derivada a memorizar para a integração
(\cref{ch:b1:integration}):
\[
\int \frac{\dd x}{1 + x^2} = \arctan x + C,
\qquad
\int \frac{\dd x}{\sqrt{1 - x^2}} = \arcsin x + C,
\qquad
\int \frac{\dd x}{\sqrt{x^2 + 1}} = \operatorname{arsinh} x + C,
\]
e, para uma quadrática geral, reduza a estes completando o quadrado
e reescalando.
\end{method}

\begin{remark}[Onde este capítulo é usado]\label{rem:b1:functions:whereused}
Este capítulo é o vocabulário de trabalho de toda a análise que vem.
A escala de crescimento da \cref{prop:b1:functions:powerrules} decide
questões de convergência ao longo de todo o \cref{ch:b1:seq,ch:b1:series}; as
fórmulas de derivada da \cref{prop:b1:functions:arcderiv} e da
\cref{prop:b1:functions:invhyp} são as primitivas mais frequentemente
necessárias no \cref{ch:b1:integration}, via o
\cref{met:b1:functions:primitive}; as \omterm{def:b1:functions:hyperbolic}{funções hiperbólicas}
parametrizam as soluções da equação $y'' = y$ no
\cref{ch:b1:diffeq} exatamente como as funções circulares parametrizam
as de $y'' = -y$. A caracterização de $\exp$ por $f' = f$,
$f(0) = 1$ (\cref{prop:b1:functions:expln}) é a semente
unidimensional da teoria das equações diferenciais lineares, e a
curva catenária $y = \cosh x$ volta entre as curvas planas do
\cref{ch:b1:curves}.
\end{remark}

\begin{remark}[Interlúdio: o programa da função inversa]\label{rem:b1:functions:interlude}
Este capítulo executou quatro vezes um mesmo programa: restringir uma
função até que ela se torne uma bijeção estritamente monótona, nomear a
inversa e transportar por ela todas as fórmulas. O que foi \emph{usado}
em cada passo --- que uma função contínua estritamente monótona num
intervalo é uma bijeção sobre um intervalo, e que a sua inversa é
contínua e depois derivável fora dos pontos críticos --- foi tomado
emprestado a crédito do ensino médio. A dívida é paga dentro deste
volume: o \cref{ch:b1:continuity} demonstra o enunciado da bijeção
(o teorema da inversa monótona, apoiado no teorema do valor
intermediário), e o \cref{ch:b1:derivative} demonstra a regra de derivação
$(f^{-1})' = 1/(f' \circ f^{-1})$, que produziu silenciosamente todas as
fórmulas da \cref{prop:b1:functions:arcderiv} e da
\cref{prop:b1:functions:invhyp}. Ler esses capítulos com
$\arcsin$ e $\operatorname{arsinh}$ em mente --- como os exemplos
trabalhados que a teoria foi construída para justificar --- é a maneira
prevista de contornar a aparente circularidade.
\end{remark}

\section{Exercícios}

\begin{exercise}[$\star$]\label{exo:b1:functions:1}
Calcule sem calculadora: $\arcsin\frac{\sqrt 3}{2}$;
$\;\arccos\bigl(-\frac 12\bigr)$; $\;\arctan(-1)$;
$\;\arcsin\bigl(\sin\frac{5\pi}{6}\bigr)$;
$\;\arccos\bigl(\cos\frac{7\pi}{4}\bigr)$.
\end{exercise}

\begin{exercise}[$\star$]\label{exo:b1:functions:2}
Dê o domínio de definição de $f(x) = \arcsin(2x - 1)$ e calcule
$f'$ onde estiver definida. Mesmas perguntas para $g(x) = \arctan\sqrt{x}$.
\end{exercise}

\begin{exercise}[$\star$]\label{exo:b1:functions:3}
Demonstre que, para todo $x \in \R$: $\cosh x + \sinh x = \eu^x$ e
$(\cosh x + \sinh x)^n = \cosh nx + \sinh nx$ para todo $n \in \Z$
(uma fórmula de De Moivre hiperbólica).
\end{exercise}

\begin{exercise}[$\star$]\label{exo:b1:functions:4}
Simplifique $\cos(2\arcsin x)$ e $\sin(2\arctan x)$, transformando-os em
expressões algébricas de $x$.
\end{exercise}

\begin{exercise}[$\star$]\label{exo:b1:functions:5}
Ordene, para $x$ grande, as funções
$x^{100}$, $\;\eu^{\sqrt x}$, $\;(\ln x)^{1000}$, $\;\eu^{x/100}$,
$\;x^{\ln x}$, do crescimento mais lento para o mais rápido, com
justificativas baseadas na \cref{prop:b1:functions:powerrules}.
\end{exercise}

\begin{exercise}[$\star\star$]\label{exo:b1:functions:6}
Estude a função $f(x) = \arctan\dfrac{2x}{1 - x^2}$ no seu domínio:
calcule $f'$, compare com $(2\arctan x)'$ e exprima $f(x)$ em
termos de $\arctan x$ em cada um dos três intervalos do domínio.
\end{exercise}

\begin{exercise}[$\star\star$]\label{exo:b1:functions:7}
Resolva em $\R$: $\;\cosh x = 2$; depois $5\cosh x - 4 \sinh x = 3$
(exprima as soluções com logaritmos). \emph{Sugestão para a segunda:
escreva tudo com $u = \eu^x$.}
\end{exercise}

\begin{exercise}[$\star\star$]\label{exo:b1:functions:8}
Demonstre a identidade $\arctan\frac{1}{2} + \arctan\frac{1}{3} =
\frac{\pi}{4}$ e depois a fórmula de Machin\index{fórmula de Machin}
\[
4\arctan\frac 15 - \arctan\frac{1}{239} = \frac{\pi}{4}.
\]
\emph{Sugestão: calcule a tangente dos dois lados usando a fórmula da
soma e controle em que intervalo cada lado se encontra.}
\end{exercise}

\begin{exercise}[$\star\star$]\label{exo:b1:functions:9}
Demonstre que, para todo $x \geq 0$: $\;x - \dfrac{x^3}{6} \leq \sin x \leq
x$ \emph{(estude as derivadas sucessivas das diferenças)} e
deduza que $\lim_{x \to 0^+} \frac{x - \sin x}{x^3} = \frac 16$ é
plausível --- o limite propriamente dito é estabelecido no
\cref{ch:b1:taylor}.
\end{exercise}

\begin{exercise}[$\star\star\star$]\label{exo:b1:functions:10}
Para $x \in \intcc{-1}{1}$, ponha
$h(x) = \arcsin\bigl(2x\sqrt{1 - x^2}\,\bigr) - 2\arcsin x$. Poderíamos
esperar $h = 0$ a partir da identidade $\sin 2y = 2\sin y\cos y$ --- mas
$h$ \emph{não} é identicamente nula. Calcule $h'$ nos intervalos abertos
em que ela existe, avalie $h$ em pontos bem escolhidos e dê a descrição
completa, constante por partes, de $h$ em $\intcc{-1}{1}$.
\end{exercise}

\begin{exercise}[$\star\star\star$]\label{exo:b1:functions:11}
(Gudermanniana) Seja $g(x) = \arctan(\sinh x)$ para $x \in \R$. Demonstre
que $g$ é uma bijeção ímpar e estritamente crescente de $\R$ sobre
$\intoo{-\frac\pi2}{\frac\pi2}$, que $g'(x) = \frac{1}{\cosh x}$,
e que $\tan g(x) = \sinh x$, $\;\sin g(x) = \tanh x$,
$\;\cos g(x) = \frac{1}{\cosh x}$: a função $g$ liga a trigonometria
circular à hiperbólica sem números complexos.
\end{exercise}

\begin{exercise}[$\star\star$]\label{exo:b1:functions:12}
Demonstre que
\[
\arctan 1 + \arctan 2 + \arctan 3 = \pi .
\]
\emph{Sugestão: calcule primeiro $\arctan 2 + \arctan 3$, localizando a
soma como no \cref{ex:b1:functions:arctansum}; atenção ao fato de que o
denominador da fórmula de adição é negativo aqui.}
\end{exercise}

\section{Problema: A equação $x^y = y^x$}

\begin{problem}\label{pb:b1:functions:1}
Que pares de números positivos satisfazem $x^y = y^x$? Todo mundo conhece
um exemplo de aparência acidental, $2^4 = 4^2 = 16$; este problema mostra
que nada nele é acidental. Toda a equação é governada
pelas variações de uma única função
\[
f(t) = \frac{\ln t}{t} \qquad (t > 0),
\]
cujo estudo fornece: o \omterm{def:b1:logic:sets}{conjunto} solução completo (uma diagonal mais um
ramo curvo passando por $(\eu, \eu)$), uma parametrização racional do
ramo, o fato de que $(2, 4)$ é o seu \emph{único} ponto inteiro
e a classificação de todos os seus pontos racionais --- além de, como
dividendos, a comparação de $\eu^\pi$ com $\pi^\eu$ e a
convergência monótona de $\bigl(1 + \frac1n\bigr)^n$ a $\eu$.
\index{função potência}

\medskip
\noindent\textbf{Parte I --- A função $f(t) = \ln t / t$.}
\begin{enumerate}
  \item Justifique que $f$ é derivável em $\intoo0{+\infty}$,
        calcule $f'$ e monte a tabela de variação: $f$
        cresce estritamente em $\intoc0\eu$, decresce estritamente em
        $\intco\eu{+\infty}$, com máximo $f(\eu) = \frac1\eu$.
  \item Determine os limites de $f$ em $0^+$ e em $+\infty$
        (\cref{prop:b1:functions:powerrules}), o sinal de $f$
        (negativa em $\intoo01$, nula em $1$, positiva além disso)
        e esboce o gráfico.
  \item Verifique por cálculo direto que $f(2) = f(4)$. (Guarde essa
        igualdade em mente: todo o problema nasce dela.)
  \item Seja $c \in \R$. Discuta, conforme o valor de $c$, o
        número de soluções de $f(t) = c$: exatamente uma para $c \leq
        0$; exatamente duas (uma em $\intoo1\eu$, outra em
        $\intoo\eu{+\infty}$) para $0 < c < \frac1\eu$; exatamente uma
        para $c = \frac1\eu$; nenhuma para $c > \frac1\eu$.
  \item Deduza que $\eu^x > x^\eu$ para todo $x > 0$ com $x \neq
        \eu$, e em particular decida qual é maior, $\eu^\pi$ ou
        $\pi^\eu$.
\end{enumerate}

\medskip
\noindent\textbf{Parte II --- A equação e a sua curva.} Nesta
parte, $x, y > 0$.
\begin{enumerate}
  \item Mostre que $x^y = y^x \iff f(x) = f(y)$.
  \item Deduza a estrutura do \omterm{def:b1:logic:sets}{conjunto} solução: todos os pares
        diagonais $(x, x)$; e os pares \emph{não triviais} ($x \neq
        y$), que satisfazem: as duas coordenadas são $> 1$, e uma
        delas está em $\intoo1\eu$, enquanto a outra está em
        $\intoo\eu{+\infty}$.
  \item Parametrize os pares não triviais: escrevendo $y = tx$ com $t
        > 0$, $t \neq 1$, mostre que $x^y = y^x$ força
        \[
        x(t) = t^{\frac1{t-1}},
        \qquad
        y(t) = t\,x(t) = t^{\frac{t}{t-1}},
        \]
        e que, reciprocamente, todo par desses é solução.
  \item Verifique que $t = 2$ dá $(2, 4)$ e demonstre a simetria
        $x(1/t) = y(t)$, $y(1/t) = x(t)$: inverter o parâmetro
        troca as duas coordenadas.
  \item Determine os limites de $x(t)$ e $y(t)$ quando $t \to 1$
        (ambos tendem a $\eu$), quando $t \to +\infty$ ($x \to 1$, $y
        \to +\infty$) e quando $t \to 0^+$ ($x \to +\infty$, $y \to
        1$). Descreva o ramo resultante: uma curva assintótica às
        retas $x = 1$ e $y = 1$, cruzando a diagonal em
        $(\eu, \eu)$.
  \item Mostre que $t \mapsto x(t)$ é estritamente decrescente em
        $\intoo1{+\infty}$. \emph{(Estude $h(t) = 1 - \frac1t - \ln
        t$, cujo sinal controla a derivada de $\ln x(t) =
        \frac{\ln t}{t - 1}$.)}
  \item Conclua que as soluções não triviais definem uma bijeção
        estritamente decrescente $\varphi \colon \intoo1\eu \to
        \intoo\eu{+\infty}$ com $x^{\varphi(x)} = \varphi(x)^x$,
        e que $\varphi$ é uma involução do ramo:
        $\varphi(\varphi(x)) = x$ sempre que os dois lados estiverem definidos.
\end{enumerate}

\medskip
\noindent\textbf{Parte III --- Pontos inteiros e racionais.}
\begin{enumerate}
  \item Demonstre que $(2, 4)$ e $(4, 2)$ são as únicas soluções
        \emph{inteiras} não triviais de $x^y = y^x$.
  \item Para $n \in \N^*$, aplique a parametrização com $t = 1 +
        \frac1n$ e mostre que
        \[
        x_n = \Bigl(1 + \frac1n\Bigr)^{n},
        \qquad
        y_n = \Bigl(1 + \frac1n\Bigr)^{n+1}
        \]
        é uma solução \emph{racional} não trivial para todo $n$,
        com $(x_1, y_1) = (2, 4)$.
  \item Reciprocamente, seja $(x, y)$ uma solução racional não trivial
        com $y > x$, e escreva $t = y/x = 1 + \frac rs$ na forma
        irredutível ($r, s \in \N^*$). Usando $x = t^{s/r}$ e admitindo
        a unicidade da fatoração em primos (familiar da
        escola; demonstrada no \cref{ch:b1:arith}), mostre que $x$
        racional força $s$ e $s + r$ a serem ambos potências $r$-ésimas
        de inteiros.
  \item Mostre, com o teorema binomial, que $a^r - b^r = r$ não tem
        soluções inteiras $a > b \geq 1$ quando $r \geq 2$, e
        conclua: as soluções racionais de $x^y = y^x$ são
        exatamente os pares $(x_n, y_n)$ da questão 14 (e as suas
        trocas).
  \item Verifique numericamente o caso $n = 2$: calcule $f(9/4)$ e
        $f(27/8)$ com quatro casas decimais e confira que coincidem.
\end{enumerate}

\medskip
\noindent\textbf{Parte IV --- Dividendos.}
\begin{enumerate}
  \item Use as questões 7 e 11 para demonstrar, sem mais
        nenhum cálculo, que a sequência $x_n = \bigl(1 +
        \frac1n\bigr)^n$ é estritamente crescente com $x_n < \eu$,
        que $y_n = \bigl(1 + \frac1n\bigr)^{n+1}$ é estritamente
        decrescente com $y_n > \eu$, e que ambas convergem para
        $\eu$. (O parâmetro $t_n = 1 + \frac1n$ decresce para
        $1$.)
  \item Estabeleça a regra geral de comparação para $1 < a < b$:
        se $\eu \leq a$, então $a^b > b^a$; se $b \leq \eu$, então
        $a^b < b^a$; e mostre, com os dois exemplos $(2, 3)$ e
        $(2, 5)$, que no caso misto $a < \eu < b$ os dois
        resultados realmente ocorrem.
  \item Suponha que a involução $\varphi$ da questão 12 seja
        derivável em $\eu$ (ela é). Derive a
        identidade $\varphi(\varphi(x)) = x$ em $x = \eu$ e deduza que
        $\varphi'(\eu) = -1$: o ramo cruza a diagonal
        perpendicularmente a ela.
  \item Encontre o único par solução com $y = 3x$, em forma
        fechada, e confira-o numericamente com quatro casas decimais
        via $f$.
  \item Entre os números $\sqrt2$, $\sqrt[3]3$, $\sqrt[4]4$,
        $\sqrt[5]5$, determine o maior e identifique os dois
        que são iguais. \emph{(Compare $n^{1/n} =
        \eu^{f(n)}$.)}
\end{enumerate}

\medskip
\noindent\textbf{Parte V --- Síntese.}
\begin{enumerate}
  \item Descreva o \omterm{def:b1:logic:sets}{conjunto} solução completo de $x^y = y^x$ no
        quadrante $x, y > 0$ --- diagonal mais ramo, a sua
        interseção $(\eu, \eu)$, as assíntotas, o ponto inteiro
        $(2, 4)$, os pontos racionais que se acumulam em
        $(\eu, \eu)$ --- numa forma que você conseguiria esboçar de
        memória.
  \item Onde exatamente o problema usou: (i) as comparações de
        crescimento da \cref{prop:b1:functions:powerrules}; (ii)
        o teorema do valor intermediário (através dos enunciados de
        bijeção); (iii) a unicidade admitida da fatoração em
        primos? Uma frase para cada.
  \item Moral, num parágrafo curto: uma única tabela de variação
        resolveu uma equação em duas incógnitas, classificou os seus
        pontos racionais e demonstrou a convergência monótona de
        $\bigl(1 + \frac1n\bigr)^n$ --- comente essa economia
        e diga onde cada fio é industrializado mais adiante no
        volume (\cref{ch:b1:seq} para a sequência,
        \cref{ch:b1:derivative} para as tabelas de variação,
        \cref{ch:b1:taylor} para a precisão que falta à tabela).
\end{enumerate}
\end{problem}
